\subsection{Why are interfaces critical in material engineering, electronics, or energy systems?}

Interfaces are not merely passive boundaries between materials; they are active, functional regions that frequently control and even dominate the physical properties of the system. As materials and devices are engineered at ever smaller scales, the role of interfaces becomes more pronounced. In conventional bulk systems, interfaces may be limited to grain boundaries or inter-phase regions. However, in nanomaterials, layered heterostructures, and device-scale systems, interfacial regions can constitute a non-negligible fraction of the total volume or surface area. Their influence is therefore not ancillary but foundational to system performance.

In materials engineering, interfaces govern a wide array of properties including adhesion, corrosion resistance, mechanical strength, and diffusion kinetics. The strength and ductility of composite systems or coatings often depend on interfacial bonding and lattice compatibility. Tailoring interfacial chemistry or morphology, through methods such as surface functionalisation, alloying, or atomic layer deposition, is often essential for achieving desired macroscopic performance metrics.

In electronic devices, interfaces critically influence charge transport, band alignment, and contact resistance. Semiconductor heterojunctions, metal-semiconductor contacts, and dielectric boundaries are fundamental to the operation of transistors, diodes, and capacitors. Localised changes in atomic registry, roughness, or defect density at the interface can significantly alter tunnelling behaviour, barrier heights, and mobility. This is especially true in 2D|3D systems, where interfacial dipoles, charge redistribution, and electrostatic reconstruction result in complex electronic responses not captured by bulk models. As devices shrink to the nanometre scale, quantum confinement and interfacial electrostatics become dominant factors, necessitating atomically precise interface control.

Energy systems are similarly interface-dominated. In lithium-ion batteries, for example, the solid-electrolyte interphase (SEI) acts as a metastable interfacial layer that dictates both stability and charge transport. The performance and longevity of solid-state batteries, fuel cells, and thermoelectric devices depend on the engineered properties of interfaces, which affect ionic conductivity, catalytic activity, thermal transport, and electronic insulation. In photovoltaics, interfacial band offsets between donor and acceptor layers, or between absorber materials and charge transport layers, govern charge separation efficiency and open-circuit voltage.

Overall, interfaces are not incidental features but active design parameters. They serve as sites of emergent phenomena, such as interfacial polarisation, bandgap renormalisation, and phonon scattering, that cannot be predicted from bulk properties alone. As such, mastering interfacial science is a gateway to developing high-performance materials and devices.

\subsection{How do crystallographic orientations, phase boundaries, or defects influence interfaces?} 

Material interfaces are regions where two or more distinct phases meet, and their atomic-level structure plays a critical role in determining the physical, electronic, and thermal properties of the overall system. These interfaces are not merely passive boundaries; rather, they are active regions that can be engineered for specific functionality, particularly in devices involving electronic, optoelectronic, and energy materials. The structural features that influence interface behaviour include crystallographic orientation, the nature of the phase boundary, and the presence of defects.

\subsubsection{Crystallographic Orientation} 

The relative orientation of crystals on either side of an interface can significantly affect the atomic registry, bond continuity, and symmetry breaking across the interface. Epitaxial alignment, where lattice planes across the interface maintain a well-defined relationship, typically results in low-strain, coherent interfaces. In contrast, a mismatch in orientation can lead to misfit dislocations or interfacial strain, influencing charge carrier transport, interfacial states, and mechanical stability. Orientation also affects the formation energy and electronic reconstruction at the interface, which can lead to phenomena like interface-induced states or altered band alignments.

\subsubsection{Phase Boundaries} 

When the interface forms between distinct structural phases, such as cubic and hexagonal polymorphs, unique bonding geometries and strain fields arise. For instance, in layered 2D|3D systems like \ce{HfS2}|\ce{HfO2}, the band alignment changes depending on whether the materials are stacked or laterally connected, with lateral interfaces showing increased reconstruction due to bonding mismatch between the different crystalline motifs. These reconstructions are not merely geometric; they result in different electrostatic and electronic behaviours, influencing band offsets and charge transfer dynamics.

\subsubsection{Defects at Interfaces} 

Defects, including vacancies, interstitials, dislocations, and grain boundaries, are often intrinsic to interface formation and can profoundly modify interfacial properties. They may serve as scattering centres, recombination sites, or sources of local doping. In grain boundaries of polycrystalline materials, for example, variations in local stoichiometry and strain can create regions of localised metallicity or enhanced dielectric response, as shown in systems exhibiting colossal permittivity. In some systems, such as CCTO (\ce{CaCu3Ti4O12}), interfacial defects and microstructural features have been shown to dominate the macroscopic dielectric properties.

In computational studies, predicting the impact of such defects requires careful treatment of their spatial distribution, charge state, and interaction with surrounding lattice atoms. Modern approaches like the RAFFLE algorithm allow for machine learning-driven structure prediction that accounts for void-filling and local bonding environments, enabling more realistic models of defect-laden interfaces.

\subsection{Examples of important technological applications involving interfaces} 

Material interfaces, where two distinct phases or materials meet, play a pivotal role in determining the performance and function of a wide array of modern technologies. These regions often exhibit structural, electronic, or thermal properties that differ markedly from those of the bulk constituents. Depending on their nature, interfaces may enhance or impede material performance. This subsection surveys several application domains in which understanding and engineering interfacial properties is essential.

In semiconductor devices, the interface between the gate dielectric and the semiconducting channel is a critical design parameter. As traditional \ce{SiO2} dielectrics are replaced by high-$\kappa$ materials such as \ce{HfO2} to meet miniaturisation demands, understanding charge transfer and band alignment across interfaces like \ce{HfS2}| \ce{HfO2} becomes essential for ensuring device reliability and performance. Notably, in such 2D|3D heterostructures, the type and strength of bonding (e.g. van der Waals vs. chemical) directly influences charge redistribution and can cause significant deviation from band alignments predicted by Anderson\rqs rule.

In layered 2D electronics, especially those using transition metal dichalcogenides (TMDCs), the stacking geometry, layer orientation, and interfacial disorder all strongly influence the band alignment. These factors affect carrier dynamics, optical transitions, and switching behaviour. Traditional models such as Anderson\rqss rule have proven insufficient; corrections using terms such as $\Delta E_\Gamma$ and $\Delta E_{\text{IF}}$ have been developed to capture hybridisation and interface dipole contributions more accurately. These corrections allow for predictive modelling of heterostructure behaviour and support targeted band offset engineering.

Interfaces also govern performance in energy storage systems. In all-solid-state batteries, for example, the interface between electrode and solid electrolyte dictates ion mobility, chemical stability, and space charge layer formation. Such interfacial processes can lead to resistive interphases or dendrite formation, reducing battery life and safety. Density Functional Theory (DFT) modelling has shown that interfacial reconstructions and electrostatic gradients can dominate device behaviour.

In thermoelectric materials, interfaces are employed to decouple electrical and thermal transport. Interface-induced phonon scattering suppresses thermal conductivity while preserving electronic conductivity; enhancing the thermoelectric figure of merit, $ZT$. Layered TMDCs and heterostructures with nanostructured interfaces exploit this mechanism to improve thermal management.

Finally, in optoelectronics and photovoltaics, interfaces play a central role in determining band bending, exciton dissociation, and charge carrier collection. In oxide-based devices, such as proposed all-oxide solar cells (e.g. \ce{CaO}|(\ce{Sn}:\ce{Ca})$_{7:1}$\ce{O}|\ce{TiO2}), interface energetics must be carefully tuned to support efficient charge separation and transport. The formation of intermediate phases or interfacial dipoles can critically affect performance, necessitating first-principles modelling for reliable device design. 

Across these examples, it is evident that interfaces are not passive boundaries, but active, tunable regions that critically determine macroscopic functionality. Their predictive modelling, through DFT, machine learning potentials, or interface-specific structure search tools like ARTEMIS, remains a central challenge in modern materials design.